\begin{tresc}
Ile jest takich liczb czterocyfrowych, w których cyfra dziesiątek jest dwa razy większa od cyfry jedności, cyfra setek jest dwa razy większa od cyfry dziesiątek, a cyfra tysięcy jest równa sumie trzech pozostałych cyfr?
\end{tresc}
\begin{wariant}{A} 1 \end{wariant}
\begin{wariant}{B} 2 \end{wariant}
\begin{wariant}{C} 3 \end{wariant}
\begin{wariant}{D} 4 \end{wariant}
\begin{wariant}{E} 5 lub więcej \end{wariant}
